% =====================================================================
%  1bat-ud1-15.tex  —  HORA 15: Taller de preparació de la prova
% =====================================================================
\horatitol{Sessió 15 — Taller de preparació de la prova}%
{Conèixer l'estructura de la prova i la rúbrica, assajar un simulacre per blocs i fer un pla d'estudi personal.}

\subsection{Idea clau}

Ens preparem per a la prova \textbf{sense sorpreses}. La prova de la sessió següent
té \textbf{quatre blocs}, un per cada criteri (C1-C4). Saber exactament què es demana
i com es corregeix fa que l'examen deixi de ser imprevisible. Després assajarem amb
un \textbf{simulacre}.

\begin{idea}
\textbf{Estructura de la prova} (un bloc per criteri):
\begin{enumerate}[leftmargin=1.6em, itemsep=2pt]
  \item \textbf{Bloc 1 (C1) — Inequacions:} resoldre una inequació de primer grau (amb
        canvi de signe) i comprovar-la.
  \item \textbf{Bloc 2 (C2) — Intervals:} escriure l'interval d'un cas real, amb els
        extrems inclosos o exclosos correctes.
  \item \textbf{Bloc 3 (C3) — Sistemes:} a partir d'un fragment de convocatòria,
        combinar dues condicions i decidir.
  \item \textbf{Bloc 4 (C4) — Notació científica:} convertir, operar i valorar
        l'ordre de magnitud.
\end{enumerate}
\textbf{Es valora}, a més del resultat: la \textbf{justificació dels passos}, l'ús
correcte de la \textbf{notació d'intervals} i la \textbf{comprovació}.
\end{idea}

\subsection{Rúbrica simplificada}

\noindent Així es corregeix cada bloc (nivells d'assoliment):
\begin{center}
\renewcommand{\arraystretch}{1.3}
\begin{tabular}{p{2.4cm} p{10.2cm}}
\toprule
\textbf{Nivell} & \textbf{Què demostra} \\
\midrule
\textbf{AE} (excel·lent) & Resol sense errors, justifica els passos i interpreta o
comprova el resultat. \\
\textbf{AN} (notable) & Resol correctament amb algun error menor; els passos i la
notació són essencialment bons. \\
\textbf{AS} (suficient) & Resol el procediment bàsic amb alguns errors; justificació
o notació incompletes. \\
\textbf{NA} (no assolit) & No identifica el procediment o comet errors que
n'invaliden el resultat. \\
\bottomrule
\end{tabular}
\end{center}

\subsection{Simulacre — un ítem de cada bloc}

\begin{exemple}[com es resol un ítem ben fet]
\textbf{Ítem tipus (Bloc 1):} resol $5-2x\ge 11$ i comprova-ho. Un \textbf{ítem
complet} mostra els passos, gira el signe i comprova:
\[
-2x\ge 6 \;\Rightarrow\; x\le -3 \quad(\text{interval } (-\infty,-3]).
\]
Comprovació amb $x=-4$: $5-2(-4)=13\ge 11$, cert. \emph{La comprovació és el que
distingeix un AE d'un AS.}
\end{exemple}

\begin{exercicis}
\textbf{Resol el simulacre amb temps controlat. Un ítem per bloc:}
\begin{exlist}
  \item \textbf{(Bloc 1 — C1)} Resol $7-3x\le 1$, dóna l'interval i comprova-la.
  \item \textbf{(Bloc 2 — C2)} Un ajut és per a joves «de $16$ anys fins a menys de
        $25$». Escriu la inequació doble i l'interval, indicant com és cada extrem.
  \item \textbf{(Bloc 3 — C3)} Una beca demana edat dins de $[18,30]$ \emph{i} renda
        $\le\num{18000}\eur$. Decideix per a un estudiant de $24$ anys amb renda
        $\num{19000}\eur$, justificant-ho.
  \item \textbf{(Bloc 4 — C4)} Un pressupost de $2,4\times 10^{8}\eur$ es reparteix
        entre $3\times 10^{5}$ persones. Calcula els euros per persona i digues si
        l'ordre de magnitud té sentit.
\end{exlist}
\end{exercicis}

\subsection{Formulari-resum (suport)}

\begin{idea}
\textbf{Tingues-ho a mà mentre prepares la prova:}
\begin{itemize}[leftmargin=1.4em, itemsep=2pt]
  \item \textbf{Regla del signe:} en multiplicar o dividir per un \textbf{negatiu},
        \textbf{gira} el símbol de la desigualtat.
  \item \textbf{Intervals:} claudàtor $[\,]$ si l'extrem hi entra ($\le$, $\ge$);
        parèntesi $(\,)$ si no ($<$, $>$); l'$\infty$ sempre amb parèntesi.
  \item \textbf{Sistemes:} cal complir \textbf{totes} les condicions; intersecció =
        part comuna.
  \item \textbf{Notació científica:} $a\times 10^{n}$ amb $1\le a<10$; multiplicar
        suma exponents, dividir resta exponents.
\end{itemize}
\end{idea}

\noindent\textbf{Pla d'estudi personal.} A partir de la teva graella d'autoavaluació
(sessió 8) i del simulacre, anota els \textbf{dos o tres tipus d'exercici} que has de
repassar amb més atenció:
\begin{center}
\renewcommand{\arraystretch}{1.5}
\begin{tabular}{c p{10.5cm}}
\toprule
\textbf{\#} & \textbf{Què he de repassar (i com)} \\
\midrule
1 & \rule{9.8cm}{0.4pt} \\
2 & \rule{9.8cm}{0.4pt} \\
3 & \rule{9.8cm}{0.4pt} \\
\bottomrule
\end{tabular}
\end{center}

% ---------------------------------------------------------------------
\fullsolucions{Sessió 15 — simulacre}
\begin{solucions}
\begin{sollist}
  \item \textbf{(Bloc 1)} $7-3x\le 1 \Rightarrow -3x\le -6$; girem: $x\ge 2$, interval
        $[2,+\infty)$. Comprovació amb $x=3$: $7-9=-2\le 1$, cert.
  \item \textbf{(Bloc 2)} «de $16$ fins a menys de $25$»: $16\le x<25$, interval
        $[16,25)$; el $16$ tancat, el $25$ obert.
  \item \textbf{(Bloc 3)} Edat $24\in[18,30]$: compleix. Renda
        $\num{19000}\le\num{18000}$? No ($\num{19000}>\num{18000}$). Com que falla una
        condició, \textbf{no} té dret a la beca.
  \item \textbf{(Bloc 4)} $\dfrac{2,4\times 10^{8}}{3\times 10^{5}}=0,8\times 10^{3}
        =8\times 10^{2}=800\eur$ per persona. Sí, té sentit (centenars d'euros).
\end{sollist}
\end{solucions}
